\section{Conclusion}
\label{Section:Conclusion}

In this work, we revisited the Dixon resultant as a practical and theoretically grounded tool for algebraic cryptanalysis. By establishing upper bounds on the Dixon matrix size that are attained by the monomial support of generic systems, we derived refined complexity estimates. Combined with an efficient implementation and a degree-aware submatrix selection strategy, our framework enables reliable elimination for polynomial systems arising in cryptographic applications.

Our results show that Dixon resultants and Gr\"obner basis methods play complementary roles in algebraic cryptanalysis. In well-determined settings, Dixon resultants are comparable to Gr\"obner bases. Gr\"obner basis methods are more effective for overdetermined systems, whereas Dixon resultants provide a structured determinant-based elimination subroutine for underdetermined settings. 

There remain several directions for future work.
\begin{enumerate}%
	\item Further analyzing AO primitives through better algebraic models for Dixon elimination, refined complexity analyses accounting for AO-specific structures, and applications to Type~III (lookup-table) primitives in combination with recent algebraic techniques~\cite{DBLP:journals/tosc/BakJGP25};

    \item Improving Dixon-resultant techniques, including rank estimation to refine Step~4 complexity estimates and optimize Step~3, the exploitation of overdetermined systems to reduce complexity, nonheuristic extraneous-factor elimination, and connections with multivariate subresultants~\cite{DBLP:journals/moc/CoxD23};
	
	\item Optimizing Dixon resultants for MQ systems, with potential applications to schemes such as UOV~\cite{DBLP:conf/eurocrypt/KipnisPG99}. Although asymptotically favorable under the compared naive Gr\"obner-basis bounds, direct Dixon elimination is generally slower than Magma in our MQ experiments and theoretically less efficient than guessing-based hybrid methods~\cite{DBLP:journals/jmc/BettaleFP09}. Improving its theoretical or practical performance on MQ systems is a direction for future work.
\end{enumerate}
